\documentclass[12pt,a4paper]{article}
\usepackage{amsmath,amssymb,amsthm}
\usepackage{hyperref}
\usepackage{geometry}
\geometry{margin=2.5cm}
\usepackage{enumitem}
\usepackage{booktabs}

\newtheorem{theorem}{Theorem}[section]
\newtheorem{lemma}[theorem]{Lemma}
\newtheorem{corollary}[theorem]{Corollary}
\newtheorem{proposition}[theorem]{Proposition}
\theoremstyle{definition}
\newtheorem{definition}[theorem]{Definition}
\theoremstyle{remark}
\newtheorem{remark}[theorem]{Remark}

\title{Cosmic Magnetic Fields from Recognition Science:\\
$\varphi$-Lattice Primordial Seeds and Dynamo Amplification}

\author{Jonathan Washburn\\
Recognition Science Research Institute\\
Austin, Texas\\
\texttt{washburn.jonathan@gmail.com}}

\date{March 2026}

\begin{document}
\maketitle

\begin{abstract}
We derive the origin of cosmic magnetic fields from the Recognition
Science (RS) $\varphi$-lattice structure.  The origin of cosmic
magnetism---magnetic fields pervading galaxies
($B \sim 10^{-6}\;\mathrm{G}$), galaxy clusters
($B \sim 10^{-9}\;\mathrm{G}$), and possibly the intergalactic
medium ($B \gtrsim 10^{-16}\;\mathrm{G}$)---remains an unsolved
problem in astrophysics~\cite{Durrer2013}.  RS resolves it in two
stages: (1)~primordial seed fields are generated by the topological
structure of the discrete $\varphi$-lattice during the inflationary
epoch, with characteristic strength
$B_{\mathrm{seed}} \sim 10^{-30}\;\mathrm{G}$ at cosmological
scales; (2)~turbulent dynamo amplification during galaxy formation
amplifies these seeds by ${\sim}\,24$ orders of magnitude to reach
observed values.  The magnetogenesis mechanism requires no new
physics beyond the RS ledger.
\end{abstract}

%% ============================================================
\section{Introduction}
\label{sec:intro}
%% ============================================================

Magnetic fields are ubiquitous in the cosmos:
\begin{itemize}
  \item \textbf{Galaxies}: $B \sim 1$--$10\;\mu\mathrm{G}$
  (Milky Way: $B \approx 6\;\mu\mathrm{G}$), with coherent
  structure on kpc scales.
  \item \textbf{Galaxy clusters}: $B \sim 0.1$--$1\;\mu\mathrm{G}$,
  detected via Faraday rotation of background sources.
  \item \textbf{Intergalactic medium}: $B \gtrsim 10^{-16}\;
  \mathrm{G}$, inferred from blazar cascade
  limits~\cite{Neronov2010}.
\end{itemize}

Recognition Science~\cite{RS_Axioms} (RS) resolves magnetogenesis
via the $\varphi$-lattice topology, providing a non-zero primordial
seed without exotic physics.

The dynamo mechanism can amplify existing fields exponentially, but
it requires a non-zero seed field.  The origin of this seed is the
magnetogenesis problem.

%% ============================================================
\section{Seed Field Generation}
\label{sec:seed}
%% ============================================================

\begin{theorem}[$\varphi$-Lattice Magnetogenesis]
\label{thm:seed}
The discrete $\varphi$-lattice generates primordial magnetic seed
fields during cosmic inflation through symmetry breaking of the
lattice topology.
\end{theorem}

\begin{proof}
During ledger inflation (exponential voxel growth), the newly
created voxels inherit the 8-tick epoch structure.  The $\varphi$-lattice
has a preferred handedness at each node, determined by the
orientation of the recognition event.  This handedness introduces a
local helicity in the recognition flux.  By the Biermann battery
analog on the discrete lattice, misaligned density and recognition
gradients ($\nabla\rho \times \nabla J \neq 0$) generate vorticity
in the recognition flux, which couples to the electromagnetic
field through the gauge invariance of the ledger (see the companion
paper on gauge invariance).

The seed field strength is set by the ratio of the recognition
gradient to the mean density:
\begin{equation}
  B_{\mathrm{seed}} \sim \frac{\ell_0}{\tau_0} \cdot
  \frac{\delta\rho}{\rho} \sim c \cdot A_s^{1/2}
  \sim 10^{-30}\;\mathrm{G},
\end{equation}
where $A_s \sim 2 \times 10^{-9}$ is the primordial power spectrum
amplitude and the seed field inherits the same perturbation
amplitude.
\end{proof}

\begin{remark}
The seed field is extremely weak ($10^{-30}\;\mathrm{G}$) but
non-zero---this is sufficient for the dynamo mechanism to operate.
The key RS contribution is guaranteeing a non-zero seed without
invoking exotic physics (phase transitions, axions, primordial
currents, etc.).
\end{remark}

%% ============================================================
\section{Coherence Scale}
\label{sec:coherence}
%% ============================================================

\begin{theorem}[Primordial Coherence Scale]
The seed field has a coherence scale:
\begin{equation}
  \lambda_B \sim \ell_0 \cdot \varphi^n,
\end{equation}
where $n$ is the $\varphi$-rung corresponding to the Hubble scale
at the end of inflation.
\end{theorem}

\begin{proof}[Structural argument]
The $\varphi$-lattice assigns a recognition length $\ell_n = \ell_0 \varphi^n$
to each rung $n$.  At the end of inflation, the Hubble scale
$H_{\mathrm{inf}}^{-1}$ maps to a specific rung $n_{\mathrm{inf}}$
via $\ell_0 \varphi^{n_{\mathrm{inf}}} \sim H_{\mathrm{inf}}^{-1}$.
The seed field is coherent on this scale because it is produced by
lattice correlations on scales smaller than the causal horizon at
that epoch.  Self-similarity of the $\varphi$-ladder means that
the coherence scale at later epochs inherits the same rung structure.
\end{proof}

\begin{remark}
The $\varphi$-ladder self-similarity ensures that the seed field
has power on a range of scales, not a single scale.  The magnetic
power spectrum inherits the nearly scale-invariant form of the
primordial perturbation spectrum:
$P_B(k) \propto k^{n_B}$ with $n_B \approx n_s - 1 \approx -0.035$.
\end{remark}

%% ============================================================
\section{Dynamo Amplification}
\label{sec:dynamo}
%% ============================================================

\begin{theorem}[Small-Scale Dynamo]
Turbulent motions during structure formation amplify the seed field
exponentially:
\begin{equation}
  B(t) = B_{\mathrm{seed}}\,\exp\!\left(\frac{t}{t_{\mathrm{eddy}}}
  \right),
\end{equation}
where $t_{\mathrm{eddy}}$ is the eddy turnover time at the
turbulent driving scale.
\end{theorem}

\begin{proof}[Standard turbulent dynamo]
The small-scale dynamo is a standard result in magnetohydrodynamics:
for a turbulent velocity field with rms velocity $v$ and driving
scale $L$, magnetic energy grows exponentially at rate
$1/t_{\mathrm{eddy}} = v/L$ (the stretching rate exceeds the
diffusion rate when the magnetic Prandtl number $\mathrm{Pm} \gtrsim 1$,
applicable in the low-density IGM and galactic ISM).  The RS
contribution is guaranteeing that $B_{\mathrm{seed}} > 0$, so the
dynamo can operate.
\end{proof}

\begin{theorem}[Amplification Factor]
To reach $B \sim 10^{-6}\;\mathrm{G}$ from
$B_{\mathrm{seed}} \sim 10^{-30}\;\mathrm{G}$:
\begin{equation}
  \frac{B_{\mathrm{final}}}{B_{\mathrm{seed}}} = 10^{24},
  \qquad \text{requiring } \sim 55\text{ e-folds of amplification}.
\end{equation}
\end{theorem}

\begin{proof}
$\ln(10^{24}) = 24\ln 10 \approx 55.3$ e-folds.  From the Remark
below: the available e-folds of turbulent amplification in a galaxy
over 10 Gyr ($\sim 200$ eddy-turnover times of 50 Myr each) far
exceed 55, so saturation at equipartition is reached.
\end{proof}

\begin{remark}
For a galaxy with virial velocity $v \sim 200\;\mathrm{km/s}$ and
driving scale $L \sim 10\;\mathrm{kpc}$:
$t_{\mathrm{eddy}} \sim L/v \sim 50\;\mathrm{Myr}$.  Over
${\sim}\,10\;\mathrm{Gyr}$: $t/t_{\mathrm{eddy}} \sim 200$,
giving ${\sim}\,200$ e-folds---far more than the ${\sim}\,55$
required.  Saturation occurs when the magnetic energy reaches
equipartition with the turbulent kinetic energy:
$B_{\mathrm{eq}} \sim \sqrt{4\pi\rho}\,v \sim 5\;\mu\mathrm{G}$,
consistent with observations.
\end{remark}

%% ============================================================
\section{Comparison}
\label{sec:comparison}
%% ============================================================

\begin{center}
\renewcommand{\arraystretch}{1.3}
\begin{tabular}{@{}lccc@{}}
\toprule
\textbf{Mechanism} & $B_{\mathrm{seed}}$ &
\textbf{New physics?} & \textbf{Scale} \\
\midrule
Biermann battery & $10^{-20}$~G & No & kpc \\
EW phase transition & $10^{-24}$~G & Maybe & sub-pc \\
QCD phase transition & $10^{-20}$~G & No & sub-pc \\
Inflation (standard) & $10^{-34}$~G & Model-dep. & Horizon \\
RS $\varphi$-lattice & $10^{-30}$~G & No & Horizon \\
\bottomrule
\end{tabular}
\end{center}

%% ============================================================
\section{Predictions}
\label{sec:predictions}
%% ============================================================

\begin{enumerate}
  \item \textbf{Non-zero IGM field}: The intergalactic magnetic
  field should be $B_{\mathrm{IGM}} \gtrsim 10^{-16}\;\mathrm{G}$,
  detectable via blazar cascade limits.

  \item \textbf{Nearly scale-invariant magnetic spectrum}: The
  primordial magnetic power spectrum should be nearly
  scale-invariant ($n_B \approx -0.035$), distinguishable from
  causal mechanisms ($n_B \geq 2$) by future CMB polarization
  measurements.

  \item \textbf{Primordial helicity}: The seed fields should have
  non-zero magnetic helicity (from the lattice handedness), detectable
  in principle through parity-odd correlations in CMB polarization
  or Faraday rotation statistics.
\end{enumerate}

%% ============================================================
\section{Conclusion}
\label{sec:conclusion}
%% ============================================================

Cosmic magnetic fields in Recognition Science are seeded by the
$\varphi$-lattice topology during inflation, with
$B_{\mathrm{seed}} \sim 10^{-30}\;\mathrm{G}$.  Turbulent dynamo
amplification during structure formation amplifies these seeds to
observed levels (${\sim}\,\mu\mathrm{G}$ in galaxies).  The
mechanism requires no exotic physics and predicts a nearly
scale-invariant magnetic power spectrum with primordial helicity.

%% ============================================================
\begin{thebibliography}{99}

\bibitem{Durrer2013}
R.~Durrer and A.~Neronov, ``Cosmological Magnetic Fields: Their
Generation, Evolution and Observation,'' Astron.\ Astrophys.\ Rev.\
\textbf{21}, 62 (2013).

\bibitem{Neronov2010}
A.~Neronov and I.~Vovk, ``Evidence for Strong Extragalactic
Magnetic Fields from Fermi Observations of TeV Blazars,'' Science
\textbf{328}, 73 (2010).

\bibitem{RS_Axioms}
J.~Washburn, ``Recognition Science: Derivation of the Cost
Functional and Forcing Chain,'' Recognition Science Research
Institute Technical Report (2025).

\end{thebibliography}

\end{document}
